\section{Uniform cross-flight summation and the complete mark count}
\label{sec:xf-pressure}
The exponential lower bound on $B_h$ does not by itself permit
multiplication by the number of physical words. We instead compare
$B_h^{-1}$ with the \emph{same} stable cocycle used in
Input~\ref{input:v112-geometric-pressure}. This comparison includes
both endpoint coordinate factors and is proved below. Only the
unconditioned pressure inequality $\sup_{R,l}P_l(R)\le Q_*$ is
used in the ensuing sum; no exact-label, roof-conditioned or new
twisted pressure estimate is assumed.

\subsection{A stable optical comparison between the two physical feet}
On the full-area, two-sided regular set $M_R'$, let $x_i$ be the
outgoing state at the $i$th intervening collision. Keep the Euclidean
arclength convention in $(s,\varphi)$, with $s=R\alpha$ and
$p=\sin\varphi$. Set
\begin{equation}\label{eq:xf-comparison-constant}
 L_\circ=\sqrt{1+(\kappa_++\tau_-^{-1})^2},\qquad
 C_J=\frac{L_\circ}{\kappa_-a_\circ^2}.
\end{equation}
These constants do not depend on $h$, $R$ or the exact labels.

\begin{lemma}[The cross-flight Green entry pays its stable weight]
\label{lem:xf-stable}
At every regular marked trajectory in $M_R'$,
\begin{equation}\label{eq:xf-stable}
 \frac1{B_h}\le C_J c_h\,j_{h-1}(x_1)
                      \le C_Jj_{h-1}(x_1),\qquad h\ge1.
\end{equation}
For $h=1$ this uses the empty stable product $j_0=1$.
\end{lemma}
\begin{proof}
We spell out the stable-line convention so that a hidden grazing
coordinate factor cannot enter the comparison. At an outgoing
collision write the stable transverse slope as $-\beta_i$.
Propagation to the next collision, followed by reflection, gives
\begin{equation}\label{eq:xf-riccati}
 \beta_i=\frac{k_{i+1}+\beta_{i+1}}
                {1+\tau(x_i)(k_{i+1}+\beta_{i+1})},
 \qquad k_i=\frac{2}{Rc_i},
 \qquad 0<\beta_i<\frac1{\tau(x_i)}\le\frac1{\tau_-}.
\end{equation}
One can obtain this invariant line without a further summability
assumption: start the backward recurrence with terminal value zero.
Every recurrence map takes $[0,\infty)$ into $[0,\tau_-^{-1}]$
and its derivative is at most
$(1+2\kappa_-\tau_-)^{-2}<1$. Differences caused by a further
terminal continuation therefore tend to zero geometrically, giving
the unique invariant negative transverse slope. It is the stable
line in the usual invariant collision cone. This agrees with the
stable-line convention in Section 1.1 of \cite{BDPressure}; no
pressure conclusion is being imported from that observation.

If $V_i=d\varphi/ds$ on this line, the identities
$J_i=c_i\,ds$ in the outgoing direction frame and
$d\vartheta=(\kappa+V_i)ds$ give
\[
 V_i=-\kappa-c_i\beta_i,\qquad
 L_i=\sqrt{1+V_i^2}\in[1,L_\circ],\qquad
 d\ell_i=|J_i|L_i/c_i.
\]
Absolute values remove the harmless alternating transverse-frame
sign. The physical stable Jacobian between collisions $1$ and $h$
therefore satisfies the \emph{exact} coordinate identity
\begin{equation}\label{eq:xf-metric}
 \frac{|J_h|}{|J_1|}
       =j_{h-1}(x_1)\frac{c_h}{c_1}\frac{L_1}{L_h}.
\end{equation}
For the first foot to the first collision, backward propagation
through that collision gives
\[
 \frac{|J_1|}{|J_{\rm first}|}
       =\frac1{1+a(k_1+\beta_1)}
       \le\frac{c_1}{2\kappa_-a_\circ}.
\]
For the final collision to the second foot the ratio is
$1-b\beta_h\in(0,1)$. The second foot precedes the next actual
collision by at least $a_\circ$; hence its stable slope magnitude
$\beta_{\rm second}$ satisfies
\begin{equation}\label{eq:xf-last-slope}
 0<\beta_{\rm second}=\frac{\beta_h}{1-b\beta_h}
     \le\frac1{\tau(x_h)-b}\le\frac1{a_\circ}.
\end{equation}
Combining these facts with \eqref{eq:xf-metric}, the stable
transverse contraction between the feet obeys
\begin{equation}\label{eq:xf-gamma}
 \gamma:=\frac{|J_{\rm second}|}{|J_{\rm first}|}
    \le \frac{L_\circ c_h}{2\kappa_-a_\circ}
                                            j_{h-1}(x_1).
\end{equation}
All incidences at the first endpoint cancel in this formula; there
is no assumption that either endpoint is uniformly nongrazing.

Since $\det\mathcal M_h=1$, applying its inverse to a vector of
slope $-\beta_{\rm second}$ gives
$\gamma^{-1}=D_h+B_h\beta_{\rm second}$. If the final free-flight
factor is removed, all remaining matrix entries are nonnegative;
thus $B_h=B_0+bD_h\ge a_\circ D_h$. Equations
\eqref{eq:xf-last-slope} and \eqref{eq:xf-gamma} now imply
\[
 \frac1{B_h}
  =\gamma\left(\frac{D_h}{B_h}+\beta_{\rm second}\right)
  \le\frac{2\gamma}{a_\circ}
  \le\frac{L_\circ c_h}{\kappa_-a_\circ^2}j_{h-1}(x_1).
\]
This proves \eqref{eq:xf-stable}. The proof uses only the positive
Jacobi propagation and the stable-line convention; the pressure
bound is not used until the next theorem.
\end{proof}

\begin{theorem}[A complete cross-flight clearance-density bound]
\label{thm:xf-pressure}
Assume the retained bound $\sup_{R,l}P_l(R)\le Q_*$ from
Input~\ref{input:v112-geometric-pressure}, with its original
physical words and stable norm. Then for every $h\ge1$ and almost
every $(u_1,u_2)\in(0,\delta_*)^2$,
\begin{equation}\label{eq:xf-root-sum}
 \sum_{a\in\mathcal A_{h,R}(u_1,u_2)}\frac1{B_h(x_a)}
                       \le4N_\circ^2C_JQ_*.
\end{equation}
Consequently the original invariant collision measure satisfies
\begin{equation}\label{eq:xf-pair-bound}
 \int N_{\delta_1}(x)N_{\delta_2}(T_R^hx)\,d\nu(x)
        \le C_{\rm cross}\delta_1\delta_2,
 \qquad C_{\rm cross}=\frac{N_\circ^2C_JQ_*}{\pi R_-},
\end{equation}
uniformly in the full separation $h\ge1$, the original radius
interval, and $0<\delta_i\le\delta_*$. Every original exact-record
or roof selector bounded by one may be retained on the left. The
essential bound in \eqref{eq:xf-root-sum} is on the \emph{two
clearance variables}, not on the physical roof.
\end{theorem}
\begin{proof}
First remove the bulk-null complement of $M_R'$ on the countably
many regular charts in Theorem~\ref{thm:xf-coarea}. Because their
two-clearance Jacobian is nonzero, the image of this complement is
null in $du_1du_2$. For each fixed $R$ and $h$, almost every pair
of clearances therefore has all its admitted roots in $M_R'$.
There are only countably many separations. A radius-dependent null
set is sufficient; no common exceptional set for all radii is
asserted.

Fix such a pair of clearances. Specify the first actual collision
center relative to the first marked center, and the second marked
center relative to the last actual collision center. This has at
most $N_\circ^2$ choices by \eqref{eq:xf-candidates}. There are
four endpoint sign pairs. For each of these choices and each
realizable physical word $w$ of length $h-1$, Theorem~\ref{thm:xf-unique}
allows at most one root. Lemma~\ref{lem:xf-stable} bounds its
weight by $C_J\sup_{[w]_R\cap M_R'}j_{h-1}$. The full positive
sum is therefore bounded by
\[
        4N_\circ^2 C_J
          \sum_{w\in\mathcal W_{h-1}(R)}
                         \sup_{[w]_R\cap M_R'}j_{h-1}
        \le4N_\circ^2 C_JQ_*.
\]
This is a weighted sum, not $D^{h-1}$ times a worst trajectory.
It is valid even when word domains are disconnected and when the
selected roots have different endpoint momenta. For $h=1$, the
sum over the empty word is exactly $P_0=1$.

Use \eqref{eq:xf-coarea} with $f=1$ and integrate the last bound
over the original clearance rectangle. The exact prefactor gives
$4/(4\pi R)$, proving \eqref{eq:xf-pair-bound}. A Borel selector
in the unchanged physical event is handled by positive domination
only after the exact coarea identity; it never becomes an
independent probability factor. The constants are independent of
$h$ because the assumed pressure bound is, not because a finite
family of tested words was extrapolated.
\end{proof}

\subsection{All chronological duplicate marks, including repeated centers}
For a finite original record of $m$ flights set
\begin{equation}\label{eq:xf-counts}
 \mathscr N_{m,\boldsymbol\delta}(x)
      =\sum_{j=0}^{m-1}N_{\delta_j}(T_R^jx),\qquad
 S_{m,\boldsymbol\delta}=\sum_{j=0}^{m-1}\delta_j,
 \quad 0<\delta_j\le\delta_*.
\end{equation}
A center seen on two different flights gives two marks. Thus this
count is different from both the number of distinct lattice centers
and the indicator of at least one witness. Let
\begin{equation}\label{eq:xf-factorial-constant}
 C_F=\max\{2C_{\rm pair},C_{\rm cross}\},\qquad
 C_{\rm pair}=\frac{N_{\rm pair}}{\pi R_-c_{\rm pair}}
\end{equation}
with the inherited same-flight constant of
Corollary~\ref{cor:dc-unique}.

\begin{theorem}[Uniform second factorial control for the complete record]
\label{thm:xf-full-count}
Under the pressure input of Theorem~\ref{thm:xf-pressure}, the
original collision measure satisfies
\begin{align}
 \int\mathscr N_{m,\boldsymbol\delta}\,d\nu
           &=S_{m,\boldsymbol\delta}/R,\label{eq:xf-first-moment}\\
 \int\mathscr N_{m,\boldsymbol\delta}
       (\mathscr N_{m,\boldsymbol\delta}-1)\,d\nu
           &\le C_F S_{m,\boldsymbol\delta}^2.\label{eq:xf-factorial}
\end{align}
The first equality needs no pressure assumption. In particular,
writing $S=S_{m,\boldsymbol\delta}$ and $\mathscr N=\mathscr N_{m,\boldsymbol\delta}$,
\begin{align}
 0\le S/R-\nu(\mathscr N\ge1)&\le\tfrac12C_FS^2,\notag\\
 \nu(\mathscr N\ge2)&\le\tfrac12C_FS^2,\notag\\
 \nu(\mathscr N=1)&\ge S/R-C_FS^2.\label{eq:xf-single-law}
\end{align}
These estimates include all pairs of collision times and all physical
words, with constants independent of $m$. They are not a Poisson
approximation and assert no independence of the marks.
\end{theorem}
\begin{proof}
Invariance of $\nu$ and the exact one-flight flux give
\eqref{eq:xf-first-moment}. Put $N_j=N_{\delta_j}(T_R^jx)$.
The exact integer identity
\[
 \mathscr N(\mathscr N-1)
      =\sum_jN_j(N_j-1)+2\sum_{j<k}N_jN_k
\]
separates same-flight ordered pairs from pairs at distinct times.
Theorem~\ref{thm:dc-factorial} bounds the first sum by
$2C_{\rm pair}\sum_j\delta_j^2$. Invariance and
Theorem~\ref{thm:xf-pressure}, applied at separation $k-j$, bound
the second by $2C_{\rm cross}\sum_{j<k}\delta_j\delta_k$.
The definition of $C_F$ proves \eqref{eq:xf-factorial}.
For a nonnegative integer $N$, use
$0\le N-\one_{\{N\ge1\}}\le N(N-1)/2$,
$\one_{\{N\ge2\}}\le N(N-1)/2$, and
$N\one_{\{N\ge2\}}\le N(N-1)$. Integration proves
\eqref{eq:xf-single-law}; no limiting interchange is needed.
\end{proof}

\begin{corollary}[The unchanged exact event and complete depth schedule]
\label{cor:xf-record}
Keep the original $\Pi=\Pi_{n,k,m,R}$, $L$, $\nu_R^*$ and all
original labels. For every Borel $0\le\psi\le1$,
\begin{align}
 0&\le\int\one_\Pi\psi(L)
             [\mathscr N_{m,\boldsymbol\delta}
                -\one_{\{\mathscr N_{m,\boldsymbol\delta}\ge1\}}]
                                              \,d\nu_R^*\notag\\
  &\le\frac{C_F}{2c_*}S_{m,\boldsymbol\delta}^2.
                                             \label{eq:xf-record}
\end{align}
For the original schedule
$\delta_j=2\varepsilon r^{\min(j,m-1-j)}$,
$r=q_0^{1/4}$, $0<2\varepsilon\le\delta_*$, this becomes
\begin{equation}\label{eq:xf-record-schedule}
 \frac{C_F}{2c_*}S_{m,\boldsymbol\delta}^2
       \le\frac{8C_F\varepsilon^2}{c_*(1-r)^2},
\end{equation}
uniformly over all $m,R,n,k$ and bounded Borel roof tests.
Equivalently, the difference between the mark-count-weighted and
one-witness-weighted exact-record roof measures is a positive
measure with this total-variation bound. This now pays
\emph{cross-flight as well as within-flight overcount}; the former
was not covered by Corollary~\ref{cor:dc-schedule}.
\end{corollary}
\begin{proof}
The integrand is nonnegative and is bounded by
$\mathscr N(\mathscr N-1)/2$. Retain the original exact event
until using $\nu_R^*=c_*^{-1}\one_{Y_R^*}\nu$ and positive
domination. Equation~\eqref{eq:xf-factorial} gives
\eqref{eq:xf-record}. There are at most two flights at each
nonnegative depth, including the middle flight when $m$ is odd.
Thus $S\le2(2\varepsilon)\sum_{d\ge0}r^d=4\varepsilon/(1-r)$,
which proves \eqref{eq:xf-record-schedule}. Taking the supremum
over Borel $0\le\psi\le1$ is exactly the total variation of the
positive difference measure. The equality of this measure with the
stated pushforward follows directly from its definition, not from
a change of conditional law.
\end{proof}

\subsection{What this theorem does, and does not, close}
The new derivative, global tangential-end uniqueness and coarea
statements are geometric theorems for the actual billiard. The
uniform complete mark-count theorem is a deduction from the
\emph{already retained} unconditioned physical pressure input, whose
independent reconstruction remains outstanding. It is not presented
as a new proof of that input. The complete double-clearance density
and the all-time duplicate payment are stronger than the inherited
same-flight assertions, but their measures are still different from
the target-roof essential height.

More precisely, for the full original exact event define the positive
roof measures
\[
 \mu^{\rm count}_{m,\boldsymbol\delta}
   =L_\#(\one_\Pi\mathscr N_{m,\boldsymbol\delta}\nu_R^*),\qquad
 \mu^{\rm one}_{m,\boldsymbol\delta}
   =L_\#(\one_\Pi\one_{\{\mathscr N_{m,\boldsymbol\delta}\ge1\}}
                                                    \nu_R^*).
\]
The exact identity
$\mu^{\rm count}=\mu^{\rm one}+\mu^{\rm excess}$ has
$\mu^{\rm excess}\ge0$ and the bound
\eqref{eq:xf-record-schedule}. The first-order one-mark source is
\emph{not} subtracted from the problem. Even in the unconditioned
law, \eqref{eq:xf-first-moment} and \eqref{eq:xf-single-law} retain
its positive first-order term $S/R$. The exact-label source need not
have that same leading coefficient; no conditioning argument here
claims otherwise.

To close the original theorem, a proof must still control the
clearance-selected first-order source and any excess which enters
the \emph{actual paired source} in the $m^2L^\infty(dt)$ topology,
including critical, degenerate, one-sided boundary and section-cut
contributions, or dominate that complete source by the original
nongrazing capacity. This bulk identity alone does not identify a
paired endpoint trace with a bulk restriction. The low-incidence
shallow source must also be paid. The inherited precise common
closure condition remains \eqref{eq:v112-remaining}; it has not
been renamed or declared true.

At every fixed outer $B$ the clearance widths are fixed first.
Neither the $O(\varepsilon^2)$ mass above nor the exponential
single-trajectory lower bound for $B_h$ is multiplied by $m^2$
and asserted to vanish. The complete physical sum, the radius and
exact-label suprema, the roof essential supremum, the collision
limsup and the final outer $B$ limit retain their original meanings
and order. Thus the raw pointwise local limit and four-coordinate
inversion target is unchanged and not yet proved.
